Let $p \in k[X_1, \ldots, \hat{X}_i, \ldots, X_{n+1}]$ be a polynomial of total degree $d$. The homogenization of $p$ at $X_i$, denoted $H_{X_i}(p)$ or $H_i(p)$, is defined by

$$H_{X_i}(p)(X_1, \ldots, X_{n+1}) = X_i^d \cdot p\!\left(\frac{X_1}{X_i}, \ldots, \frac{\widehat{X_i}}{X_i}, \ldots, \frac{X_{n+1}}{X_i}\right).$$

The resulting polynomial is homogeneous of degree $d$. When the variable to be homogenized at is clear from context, we write $H(p)$.
